\section{Remove Duplicates from a Sorted DLL}
\label{sec:ll-remove-duplicates-dll}

\problemstatement{Given a doubly linked list sorted in ascending order,
remove nodes with duplicate values so each value appears once, returning the
head.}

\begin{patternbox}
In a sorted list, duplicates are \textbf{adjacent}, so a single pass
comparing each node to the next -- and splicing out equal successors --
removes all duplicates. The \texttt{prev} pointer keeps the deletions
$O(1)$.
\end{patternbox}

\subsection*{Skip equal neighbours}

Walk from the head. While the current node's value equals its successor's,
unlink the successor (rewire the current node's \texttt{next} and the
node-after's \texttt{prev}). Only advance to the next distinct value once no
more equal successors remain. Sortedness guarantees all copies of a value
are consecutive, so this local check suffices.

\begin{ideabox}[Sorted Means Duplicates Are Adjacent]
Because equal values sit next to each other in a sorted list, comparing
only neighbours catches every duplicate -- no hashing or full scan needed.
The DLL links make each removal constant time.
\end{ideabox}

\subsection*{Algorithm}

\begin{enumerate}
  \item \texttt{cur = head}.
  \item While \texttt{cur} and \texttt{cur->next}: if equal, splice out
        \texttt{cur->next}; else advance \texttt{cur}.
  \item Return the head.
\end{enumerate}

\subsection*{Dry run}

\dryrun{Sorted DLL $1,1,2,3,3$.}

\begin{itemize}
  \item At first $1$: remove the second $1$. At $2$: no dup. At first $3$:
        remove the second $3$.
  \item Result $1,2,3$.
\end{itemize}

\subsection*{C++ implementation}

\begin{lstlisting}
#include <bits/stdc++.h>
using namespace std;

struct Node {
    int val; Node* prev; Node* next;
    Node(int v) : val(v), prev(nullptr), next(nullptr) {}
};

Node* build(const vector<int>& a) {
    Node dummy(0), *t = &dummy;
    for (int x : a) { Node* n = new Node(x); t->next = n; n->prev = t; t = n; }
    Node* head = dummy.next; if (head) head->prev = nullptr;
    return head;
}

Node* removeDuplicates(Node* head) {
    Node* cur = head;
    while (cur && cur->next) {
        if (cur->val == cur->next->val) {
            Node* dup = cur->next;
            cur->next = dup->next;                 // splice out duplicate
            if (dup->next) dup->next->prev = cur;
        } else {
            cur = cur->next;
        }
    }
    return head;
}

void print(Node* h) { for (; h; h = h->next) cout << h->val << " "; cout << "\n"; }

int main() {
    print(removeDuplicates(build({1, 1, 2, 3, 3})));   // 1 2 3
    return 0;
}
\end{lstlisting}

\begin{complexitybox}
\textbf{Time Complexity.} $O(n)$ -- one pass. \textbf{Space Complexity.}
$O(1)$.
\end{complexitybox}

\begin{takeawaybox}
On a sorted list, duplicates are neighbours, so a neighbour-comparison pass
removes them all. Recognising the sorted invariant turns a potential
hashing problem into an $O(1)$-space local sweep.
\end{takeawaybox}
